% Part: proof-theory
% Chapter: natural-deduction
% Section: translation-G2i

\documentclass[../../../include/open-logic-section]{subfiles}

\begin{document}

\olfileid{pt}{nat}{tgi}

\olsection{\Log{G2i} पासून \Log{N2i} मध्ये रूपांतर}

~\Log{G2i} मधील क्रमवर्तींचे पूर्वपक्ष हे
!!{formula}s यांचे बहुसंच~$\Gamma$ असतात; तर
\Log{N2i} मधील क्रमवर्तींचे पूर्वपक्ष हे
खूण लावलेल्या सूत्रांचे संच~$\Gamma'$ असतात.
त्यांतील प्रत्येक खूण फक्त एकदाच येते.
खूण लावलेल्या !!{formula}s चा संच~$\Gamma'$
हा बहुसंच~$\Gamma$ शी \emph{अनुरूप आहे} असे
तेव्हाच म्हणू, जेव्हा प्रत्येक !!{formula}~$!A$
साठी $x : !A$ या रूपातील~$\Gamma'$ चे
!!{element}s जितके आहेत तितकी संख्या
~$\Gamma$ मध्ये~$!A$ जितक्या वेळा येते
(म्हणजे~$!A$ च्या प्रती जितक्या आहेत) त्या
संख्येपेक्षा अधिक नसते.

\begin{prop}\ollabel{prop:G2i-to-N2i}
$\Log{G2i} + \Cut \Proves \Gamma \Sequent \Delta$
असेल, तर
$\Log{N2i} \Proves \Gamma' \Sequent !C$.
येथे~$\Delta$ रिकामा असेल तर
$!C \ident \lfalse$; अन्यथा
$\Delta = \{!A\}$ आणि $!C \ident !A$.
तसेच~$\Gamma'$ हा~$\Gamma$ शी अनुरूप आहे.
\end{prop}

\begin{proof}
  ~$\pi$ ही~\Log{G2i} + \Cut मध्ये
  $\Gamma \Sequent \Delta$ ची
  !!a{proof} असेल, तर~$\Gamma$
  शी अनुरूप असा~$\Gamma'$ आणि
  ~\Log{N2i} मध्ये
  ~$\Gamma' \Sequent !C$ ची
  !!a{proof}~$\delta$ मिळते,
  हे दाखवू. त्यासाठी
  $n = \pheight{\pi}$
  यावर विगमन करू.

  $n = 0$ असेल, तर~$\pi$
  हे स्वयंसिद्धक आहे.
  ते $\lfalse \Sequent \quad$
  असे असेल, तर~$\delta$
  हे $x:\lfalse \Sequent \lfalse$
  स्वयंसिद्धक; म्हणजे
  $\Gamma' = \{x:\lfalse\}$
  आणि $!C \ident \lfalse$.
  ~$\pi$ हे $!A \Sequent !A$
  या रूपाचे स्वयंसिद्धक असेल,
  तर~$\delta$ हे
  $x: !A \Sequent !A$.

  $n>0$ असेल, तर~$\pi$
  चा शेवट~$R$ या नियमात होतो.
  विगमन गृहीतकानुसार त्या नियमाच्या
  आधारविधानांशी अनुरूप क्रमवर्तींच्या
  !!{proof}s ~\Log{N2i} मध्ये मिळतात.
  या सिद्धतांतील सर्व मुक्तीच्या खुणा
  वेगवेगळ्या आहेत, आणि एखादी खूण~$x$
  मुक्त होत असेल तर तिला मुक्त करणाऱ्या
  अनुमानाच्या वरच येते, असे गृहीत धरू
  शकतो; गरज पडल्यास खुणा बदलू शकतो.
  या !!{proof}s एकत्र करून योग्य
  $\Gamma' \Sequent !C$ ची सिद्धता
  कशी मिळते ते आता~$R$ नुसार पाहू.

  \begin{enumerate}
    \item $R$ हा \RightR{\Weakening} असेल,
    तर~$\pi$ चा शेवट असा होतो:
    \[
    \AxiomC{}
    \RightLabel{$\pi_1$}
    \Deduce$\Gamma \fCenter$
    \RightLabel{\RightR{\Weakening}}
    \UnaryInf$\Gamma \fCenter !A$
    \DisplayProof
    \]
    ~$\pi_1$ वर विगमन गृहीतक लावल्यास
    $\Gamma' \Sequent \lfalse$ ची
    !!a{proof} मिळते. तिला~$\FalseInt$
    लावून $\Gamma' \Sequent !A$ ची
    !!a{proof} मिळते.
    \item $R$ हा \LeftR{\Weakening} असेल,
    तर~$\pi$ चा शेवट असा होतो:
    \[
    \AxiomC{}
    \RightLabel{$\pi_1$}
    \Deduce$\Gamma \fCenter \Delta$
    \RightLabel{\LeftR{\Weakening}}
    \UnaryInf$!A, \Gamma \fCenter \Delta$
    \DisplayProof
    \]
    ~$\pi_1$ वर विगमन गृहीतक लावल्यास
    $\Gamma' \Sequent !C$ ची
    !!a{proof} मिळते; तीच~$\delta$ घेऊ.
    ~$\Gamma'$ हा~$\Gamma$ शी अनुरूप असेल,
    तर तो $!A, \Gamma$ शीही अनुरूप असतो.
    कारण~$\Gamma'$ मधील $x : !A$ ची संख्या
    ही~$\Gamma$ मधील~$!A$ च्या प्रतींपेक्षा
    अधिक नाही; त्यामुळे ती
    $!A, \Gamma$ मधील~$!A$
    च्या प्रतींपेक्षाही अधिक नाही.
    \item $R$ हा \LeftR{\Contraction} असेल,
    तर~$\pi$ चा शेवट असा होतो:
    \[
    \AxiomC{}
    \RightLabel{$\pi_1$}
    \Deduce$!A, !A, \Gamma \fCenter \Delta$
    \RightLabel{\LeftR{\Contraction}}
    \UnaryInf$!A, \Gamma \fCenter \Delta$
    \DisplayProof
    \]
    ~$\pi_1$ वर विगमन गृहीतक लावल्यास
    $\Pi, \Gamma' \Sequent !C$ ची
    !!a{proof}~$\delta_1$ मिळते;
    येथे $\Pi \subseteq \{x : !A,
    y:!A\}$. $\Pi = \{x : !A, y:!A\}$
    असेल, तर~$\delta_1$ मधील
    सर्व खूण लावलेली !!{formula}s~$y:!A$
    ही~$x:!A$ ने बदला.
    ~$\delta_1$ च्या अंतिम क्रमवर्तीच्या
    संदर्भात~$x$ असल्यामुळे
    ~$\delta_1$ मधील कोणतेही अनुमान
    $x:!B$ या रूपाचे !!a{formula}
    मुक्त करत नाही, असे धरता येते;
    म्हणजे~$\delta_1$ मधील
    कोणत्याही क्रमवर्तीच्या डाव्या बाजूवरील
    $x:!B$ सक्रिय नसते.
    म्हणून बदललेली सिद्धता अजूनही
    ~\Log{N2i} मधील !!a{proof} असते.
    \item $R$ हा $\RightR{\lif}$ असेल,
    तर~$\pi$ चा शेवट असा होतो:
    \[
    \AxiomC{}
    \RightLabel{$\pi_1$}
    \Deduce$!A, \Gamma \fCenter !B$
    \RightLabel{\RightR{\lif}}
    \UnaryInf$\Gamma \fCenter !A \lif !B$
    \DisplayProof
    \]
    विगमन गृहीतकानुसार
    $x: !A, \Gamma' \Sequent !B$
    किंवा $\Gamma' \Sequent !B$
    ची !!a{proof}~$\delta_1$ मिळते;
    येथे~$\Gamma'$ हा~$\Gamma$ शी
    अनुरूप आहे. ~$\delta_1$ ला
    \Intro{\lif} चे एक अनुमान जोडून
    मिळणारी सिद्धता~$\delta$ घेऊ.
    \item $R$ हा \LeftR{\lif} असेल,
    तर~$\pi$ चा शेवट असा होतो:
    \[
    \AxiomC{}
    \RightLabel{$\pi_1$}
    \Deduce$\Gamma_1 \fCenter !A$
    \AxiomC{}
    \RightLabel{$\pi_2$}
    \Deduce$!B, \Gamma_2 \fCenter \Delta$
    \RightLabel{\LeftR{\lif}}
    \BinaryInf$!A \lif !B, \Gamma_1, \Gamma_2 \fCenter \Delta$
    \DisplayProof
    \]
    विगमन गृहीतकानुसार !!{proof}s मिळतात:
    $\delta_1$ ही $\Gamma_1' \Sequent !A$ ची;
    आणि~$\delta_2$ ही
    $x: !B, \Gamma_2' \Sequent !C$
    किंवा $\Gamma_2' \Sequent !C$ ची.
    ~$\delta_2$ दुसऱ्या रूपात असेल,
    तर~$\delta_2$ हीच~$\delta$ घेऊ.
    अन्यथा~$\delta_1$ मधील
    सर्व सूत्रांच्या आणि मुक्तीच्या खुणा
    बदलून~$\delta_1$ व~$\delta_2$
    यांत कोणतीही खूण समान नसल्याची
    खात्री करू शकतो. मग
    $\Gamma_1' \cup \Gamma_2'$ हा
    $\Gamma_1 \cup \Gamma_2$ शी अनुरूप आहे.
    ~$\delta_2$ च्या अंतिम क्रमवर्तीत
    $x:!B$ असल्यामुळे
    ~$\delta_2$ मधील~$x$ या
    मुक्ती-खुणेच्या कोणत्याही अनुमानात
    $x:!B$ सक्रिय नसते.
    ~$\delta_2$ मधील प्रत्येक
    $x:!B \Sequent !B$ स्वयंसिद्धकाच्या
    जागी पुढील !!{proof}~$\delta_3$ ठेवा:
    \[
    \Axiom$x: !A \lif !B \fCenter !A \lif !B$
    \AxiomC{}
    \RightLabel{$\delta_1$}
    \Deduce$\Gamma_1' \fCenter !A$
    \RightLabel{\Elim{\lif}}
    \BinaryInf$x: !A \lif !B, \Gamma_1' \fCenter !B$
    \DisplayProof
    \]
    आणि~$\delta_2$ मधील प्रत्येक
    $x:!B$ च्या जागी $x:!A \lif !B$
    ठेवा. परिणामी
    $\Subst{\delta_2}{\delta_3}{x:!B}$
    ही $x:!A \lif !B, \Gamma_1', \Gamma_2' \Sequent !C$
    ची !!a{proof} आहे.
    \item $R$ हा \RightR{\land} असेल,
    तर~$\pi$ चा शेवट असा होतो:
    \[
    \AxiomC{}
    \RightLabel{$\pi_1$}
    \Deduce$\Gamma_1 \fCenter !A$
    \AxiomC{}
    \RightLabel{$\pi_2$}
    \Deduce$\Gamma_2 \fCenter !B$
    \RightLabel{\RightR{\land}}
    \BinaryInf$\Gamma_1, \Gamma_2 \fCenter !A \land !B$
    \DisplayProof
    \]
    विगमन गृहीतकानुसार
    $\Gamma_1' \Sequent !A$ ची
    सिद्धता~$\delta_1$ आणि
    $\Gamma_2' \Sequent !B$ ची
    सिद्धता~$\delta_2$ मिळते.
    ~$\delta_2$ मधील सर्व सूत्रांच्या
    आणि मुक्तीच्या खुणा बदलून
    ~$\delta_1$ व~$\delta_2$ यांत
    कोणतीही खूण समान नसल्याची
    खात्री करता येते. मग
    $\Gamma_1' \cup \Gamma_2'$ हा
    $\Gamma_1 \cup \Gamma_2$ शी अनुरूप आहे.
    ~$\delta_1$ आणि~$\delta_2$
    यांच्या अंतिम क्रमवर्तींना
    $\Intro{\land}$ लावून
    !!a{proof}~$\delta$ मिळते.
    \item $R$ हा $\LeftR{\land}$ असेल,
    तर, उदाहरणार्थ, ~$\pi$ चा
    शेवट असा होतो:
    \[
    \AxiomC{}
    \RightLabel{$\pi_1$}
    \Deduce$!A, \Gamma \fCenter \Delta$
    \RightLabel{\LeftR{\land}}
    \UnaryInf$!A \land !B, \Gamma \fCenter \Delta$
    \DisplayProof
    \]
    विगमन गृहीतकानुसार
    $x: !A, \Gamma' \Sequent !C$
    किंवा $\Gamma' \Sequent !C$
    ची !!a{proof}~$\delta_1$ मिळते;
    येथे~$\Gamma'$ हा~$\Gamma$ शी
    अनुरूप आहे. दुसऱ्या प्रसंगी
    ~$\delta_1$ हीच~$\delta$ घेऊ.
    पहिल्या प्रसंगी~$x:!A$
    अंतिम क्रमवर्तीत असल्यामुळे
    ~$\delta_1$ मधील~$x$ या
    मुक्ती-खुणेच्या कोणत्याही अनुमानात
    $x:!A$ सक्रिय नसते.
    प्रत्येक $x:!A \Sequent !A$
    स्वयंसिद्धकाच्या जागी
    पुढील !!{proof}~$\delta_2$ ठेवा:
    \[
    \Axiom$x: !A \land !B \fCenter !A \land !B$
    \RightLabel{\Elim{\land}}
    \UnaryInf$x: !A \land !B \fCenter !A$
    \DisplayProof
    \]
    आणि प्रत्येक~$x:!A$ च्या जागी
    $x:!A \land !B$ ठेवा; म्हणजे
    $\Subst{\delta_1}{\delta_2}{x:!A}$
    घ्या. ही $x:!A \land !B, \Gamma' \Sequent !C$
    ची !!a{proof} आहे.
    \item $R$ हा \Cut असेल,
    तर~$\pi$ चा शेवट असा होतो:
    \[
    \AxiomC{}
    \RightLabel{$\pi_1$}
    \Deduce$\Gamma_1 \fCenter !A$
    \AxiomC{}
    \RightLabel{$\pi_2$}
    \Deduce$!A, \Gamma_2 \fCenter \Delta$
    \RightLabel{\Cut}
    \BinaryInf$\Gamma_1, \Gamma_2 \fCenter \Delta$
    \DisplayProof
    \]
    विगमन गृहीतकानुसार
    $\Gamma_1' \Sequent !A$ ची
    सिद्धता~$\delta_1$ आणि
    $x:!A, \Gamma_2' \Sequent !C$
    किंवा $\Gamma_2' \Sequent !C$
    ची सिद्धता~$\delta_2$ मिळते.
    दुसऱ्या प्रसंगी~$\delta$
    म्हणून~$\delta_2$ घ्या.
    अन्यथा, गरज असल्यास
    ~$\delta_1$ मधील खुणा बदलून
    दोन सिद्धतांच्या खुणा वेगळ्या करा.
    ~$\delta_2$ मधील प्रत्येक
    $x:!A \Sequent !A$ स्वयंसिद्धकाच्या
    जागी~$\delta_1$ ठेवून
    $\Subst{\delta_2}{\delta_1}{x:!A}$
    ही $\Gamma_1', \Gamma_2' \Sequent !C$
    ची !!{proof}~$\delta$ मिळते.
  \end{enumerate}
\end{proof}

\end{document}
